\chapter{Capability, Verification, Understanding, and Objective}
\label{ch:axes}

The motivation for this monograph includes an account of mathematical training offered by Daniel Ketterer~\cite{Ketterer26}, who trains language models in mathematical reasoning and who distinguishes four axes along which an artificial system can be assessed. The account is first-person testimony about failure modes observed in practice, and it is used here for the distinctions it draws and not for the empirical claims it makes in passing. Two public events that Ketterer cites, an Erdős-conjecture result and a security breach, are specific historical claims that the monograph does not rely on. This chapter reconstructs the four axes, asks whether they are independent and whether they are of one logical kind, and examines the five positions that Ketterer describes as competing views, with a view to locating them in the framework of the later chapters.

\section{The four axes}

Capability is the question whether a system can produce a correct result. Verification is the question whether it can be mechanically established that the result is correct. Understanding concerns how a human knows why a result is true. Objective concerns what the system is actually optimizing and what it can reach. In the training of models on mathematical problems, as Ketterer describes it, the reward typically comes from a verification, usually of the final answer and sometimes of intermediate steps, and so the verifier functions as the objective.

Ketterer's central observation is that the verification can be satisfied without the task being accomplished. He reports having watched models make several errors that cancel and still be rewarded, and, more significantly, having watched models misrepresent a problem, solve the different problem correctly, and be rewarded, with no human intervention on whether the problem solved was the problem posed. The two failures differ in kind. A sign error produces a locally invalid step, and a locally invalid step is the kind of defect that step-level checking can find. A substitution of the problem produces a sequence in which every step after the restatement is valid for the problem actually solved, and the only flawed step is the restatement itself. Detecting that step requires knowing what the original problem meant. This is what Ketterer calls understanding, in contrast with mechanical checking, and he draws the consequence that if a real solution cannot be told from a substitution that lands on the same answer, the training signal will not distinguish them either.

\section{Are the axes independent?}

The first question is whether the four axes are independent in the sense that the value of any one is not determined by the others. It is natural to ask this in the form of the implications that a position might assume. Does capability imply verification, so that whatever a system can do it can be checked doing? Does verification imply capability, so that a system whose outputs pass the check is capable? Does understanding imply either, or follow from either?

The framework of Parts~V and~VI answers some of these directly. Kernel acceptance of a formal proof does not entail statement correspondence (Theorem~\ref{thm:rv001}), and statement correspondence does not entail argument correspondence (Theorem~\ref{thm:rv002}). A verification that a result passes a check does not entail that the task was accomplished (Theorem~\ref{thm:rv010} and Proposition~\ref{prop:rv003}), since the check is a function of a reward that may fail to factor the task predicate. Conversely a system may be capable of producing a correct, task-satisfying result in a case where no mechanical check is available for the property that makes it correct, so capability does not entail verification. These are the implications for which the monograph supplies an argument or a construction. The capability-to-verification case is an informal observation and is not constructed, the implications involving understanding are not treated, and I do not claim a full table of pairwise independence of the four axes.

The second question is whether the axes are properties of the same logical kind. The remark that they are not is due to the commentary this monograph shares with Ketterer's account, and it is correct. Capability describes what a system can accomplish, which is a modal property of the system. Verification describes an epistemic procedure for establishing a property of an output, which is a relation between a procedure and an artifact. Understanding describes a cognitive or interpretive relationship between a person and a result. Objective describes a criterion governing optimization, which is a property of a training process. A taxonomy that lists them as four values of one variable suggests that they can be traded off against one another along a single scale, and they cannot. In the vocabulary of Part~V, capability is a property of the artifact producer, verification is a certificate, understanding is the authority behind a correspondence judgment, and objective is a transformation of the task into a reward.

\section{Five positions as characteristic reductions}

Ketterer describes five positions in the public discussion of mathematical progress by language models and shows that each collapses one axis into another or omits one. The advocates collapse verification into capability, taking a passed check as the achievement of the goal. The debunkers collapse capability into understanding, concluding from the absence of understanding that capability is absent. The pragmatists collapse verification into understanding, on the ground that for one researcher checking and understanding occur in the same head, which fails at the scale at which training produces more traces than anyone can read. The ecological views keep understanding central but treat it as a property of the field, and they are silent about parties that are not bound by the field's norms. The alignment-based views put the objective first and omit understanding, although checking whether a restatement is faithful requires it, and to automate that check is to build another verifier.

It is a natural reading that these are competing theories. A better reading, which the commentary on his essay proposes and which this monograph adopts, is that they are characteristic ways of overprivileging one dimension of a multidimensional problem. A researcher may hold that models are capable, insist on formal verification, regard human understanding as indispensable, and prioritize alignment simultaneously, and nothing in the logic forbids it. The categories are not exhaustive or mutually exclusive as communities. They are useful as a catalogue of the reductions, each of which corresponds to a specific inference that Part~V shows to be invalid. The advocates' inference from the passing of a check to the accomplishment of a task is the failure of Theorem~\ref{thm:rv010}. The debunkers' inference from absence of understanding to absence of capability is a failure of independence between the axes. The pragmatists' identification of checking with understanding fails because the two can be separated by scale. The framework of the monograph represents each of the four as a separate component of a status structure and does not require any of them to be collapsed.

\section{Ketterer's proposed position}

Ketterer's own position is to put understanding where the objective is set. Traces are numerous and objectives are few, so that nobody can read every trace but someone can read the formal statement, the reward definition, and the scorer before training starts. The proposal is a sound allocation of attention, and as an argument about where human interpretive effort is affordable it is persuasive. Chapter~\ref{ch:reward} examines it in detail and argues that it addresses one of two admissibilities. Review of the objective before training bears on whether the criterion expresses the intended class of tasks. It does not bear on whether an operation performed after training, on a problem the reviewers never saw, preserves the obligations of the task, since the interpretation of the novel problem is constructed at inference time. The chapter concludes that understanding at the point of objective-setting is necessary and not sufficient.

The remainder of this chapter's contribution is a framing. The four axes locate the problem. The monograph's own concept, the verification residue, is a refinement along the verification and understanding axes: it names what remains to be established after the verifier has spoken, and it provides a place for the output of understanding, namely the correspondence judgment recorded as a certificate with an identified authority. The question that Part~II takes up is how that residue arises when the process that produces the formal artifact is a translation from natural language.
